\documentclass[11pt]{article}

\usepackage[utf8]{inputenc}
\usepackage[T1]{fontenc}
\usepackage{amsmath,amssymb}
\usepackage{booktabs}
\usepackage{graphicx}
\usepackage{hyperref}
\usepackage[margin=1in]{geometry}
\usepackage{url}

\hypersetup{colorlinks=true, linkcolor=blue, citecolor=blue, urlcolor=blue}

\title{PSF-Zero: An Analytic Two-Qubit Gate Synthesizer Combined with\\
a Verified, Ordering-Aware Layout Search for Quantum Circuit\\
Transpilation}

\author{TN Holdings\\
\texttt{love.os.architect@proton.me}\\
\texttt{https://github.com/TN-Holdings-LLC/psf-zero}}

\date{\today}

\begin{document}
\maketitle

\begin{abstract}
We present PSF-Zero, a two-part quantum circuit compiler consisting of
(1) an analytic two-qubit gate synthesizer based on the Cartan (KAK)
decomposition, implemented as a compiled Rust core, and (2) a qubit
layout search built around a specific ordering strategy shown effective
in a companion paper. On a fixed 15-qubit, dense-pair-block circuit
family, the gate synthesizer reproduces Qiskit's own output to a
fidelity of $1.000000000000$ at small scale and achieves a
$7.0$--$9.3\times$ cumulative speedup over Qiskit's default
\texttt{optimization\_level=3} synthesis path across independently
repeated 10,000-iteration measurements; we report the variance of this
speedup honestly, including a $34$--$54\times$ reduction in timing
variance observed between two measurement sessions whose cause we
investigated with the actual source diff from both sessions and could
not identify. The layout search, verified through the public
\texttt{compile\_for\_hardware()} entry point rather than as an isolated
component, finds a valid embedding on 26 of 26 saturated test
instances---up from 0 of 26 before a feasibility-check defect we found
and fixed---with all 26 outputs independently verified structurally
valid (zero coupling-map violations) and, at small scale,
unitary-equivalent to the input circuit to machine precision (minimum
fidelity $0.999999999999994$ across 36 independent checks). We further
verify the same layout search on two topologies distinct from the
square grids used in development, including one (\texttt{brick}) on
which the search's own earlier ordering strategies had previously found
nothing at all. We report every defect found in the course of this
verification, including one whose root cause we mischaracterized in an
earlier stage of this investigation and later corrected, and one
inconsistency in our own reported timing figures that we found and
correct in this paper. All code, raw data, and the complete
chronological verification record are public.
\end{abstract}

\section{Introduction}

Quantum circuit compilation for near-term hardware combines two largely
independent problems: synthesizing an arbitrary computation into a
device's native gate set, and mapping the circuit's logical qubits onto
physical ones subject to connectivity constraints. PSF-Zero addresses
both. Its gate synthesis path replaces Qiskit's own general-purpose
two-qubit decomposition with a specialized, analytic KAK/Cartan
decomposition compiled to native code. Its layout path is a prototype
search built around a specific empirical finding---that a particular,
extremely simple node-ordering strategy resolves a class of instances
on which Qiskit's own production layout pass fails outright---reported
in detail in a companion paper~\cite{vf2cliff}.

This paper's contribution is not the ordering insight itself, which
belongs to the companion paper; it is the engineering and verification
of a system that uses it. We treat "the search finds a layout" as
necessary but insufficient, and verify three further properties that,
to our knowledge, are not jointly established for any comparable
research prototype: that every found layout respects the device's
connectivity (structural validity), that the resulting circuit computes
the same operator as the input (unitary equivalence, checked exactly at
small scale rather than assumed), and that the underlying search
strategy is not overfit to the square-grid topology it was developed
on.

We are equally direct about a result we cannot explain. Over the course
of this work we observed the gate synthesizer's own timing variance
drop by more than an order of magnitude between two measurement
sessions using nominally the same benchmark. We obtained the actual
source diff between the two sessions' codebases, traced every
candidate mechanism we could identify, ruled out each one by direct
test, and did not find the cause. We report this as a closed,
unresolved question rather than an unsupported explanation, consistent
with the verification standard we apply to every positive claim in this
paper.

\subsection{Contributions}

\begin{enumerate}
\item An analytic two-qubit gate synthesizer achieving
$7.0$--$9.3\times$ cumulative speedup over Qiskit's default synthesis
path, with exact correctness verified at small scale
(Section~\ref{sec:synth}).
\item An engineered layout search, built on the ordering finding of a
companion paper, verified end-to-end through the system's own public
API: 26/26 instance coverage (up from 0/26 before a defect fix),
zero structural violations, and unitary equivalence to machine
precision (Section~\ref{sec:layout}).
\item Generalization evidence on two topologies outside the search's
own development set, including a topology on which its predecessor
strategies previously found nothing (Section~\ref{sec:generalization}).
\item A full account of defects found during development and
verification, including one root-cause mischaracterization we
corrected mid-investigation and one figure-transcription error we
correct in this paper (Section~\ref{sec:defects}).
\item An honestly unresolved empirical finding---the gate synthesizer's
own variance reduction across sessions---reported with the
investigation that failed to explain it (Section~\ref{sec:variance}).
\end{enumerate}

\section{Gate synthesis}
\label{sec:synth}

\subsection{Design}

PSF-Zero's \texttt{compile()} entry point collects two-qubit blocks from
a circuit (via Qiskit's own \texttt{Collect2qBlocks} and
\texttt{ConsolidateBlocks} passes, reused rather than reimplemented) and
replaces each consolidated block's unitary with an explicit KAK/Cartan
decomposition: two single-qubit unitaries, an entangling core expressed
in the target basis (\texttt{cx} in all measurements reported here), and
two further single-qubit unitaries. The decomposition itself runs in a
compiled Rust extension (\texttt{psf\_zero\_core}), called once per
consolidated block.

\subsection{Correctness}

At small scale ($n=6$), where exact operator comparison is tractable,
the synthesized circuit reproduces the input circuit's own unitary to a
fidelity of $1.000000000000$ across all three arms compared (Qiskit
\texttt{optimization\_level=3}, PSF-Zero with and without its own
verification pass enabled). We do not compute exact unitary equivalence
at the 15-qubit scale used for the timing measurements below; this is
the same limitation stated plainly in our own project documentation,
and we repeat it here rather than let the small-scale result imply
more than it establishes.

\subsection{Performance}

\begin{table}[h]
\centering
\begin{tabular}{lrrrr}
\toprule
Session & Total (Qiskit) & Total (PSF-Zero) & Speedup & Std.\ dev.\ (PSF-Zero) \\
\midrule
Archive (4 days prior) & $92.07$\,s & $13.14$\,s & $7.00\times$ & $0.902$\,ms \\
Today & $103.75$\,s & $11.86$\,s & $8.75\times$ & $0.123$\,ms \\
Repeat & $106.56$\,s & $11.49$\,s & $9.27\times$ & $0.164$\,ms \\
\bottomrule
\end{tabular}
\caption{Cumulative time over $10{,}000$ iterations, 15-qubit
dense-pair-block circuits, \texttt{verify=False}, three independent
sessions. Standard deviation is per-call, over the same $10{,}000$
calls. \emph{The archive session's speedup was previously reported as
$8.87\times$ in our own project record; recomputed directly from the
session's own raw per-call data for this paper, it is $7.00\times$. We
correct the earlier figure here (Section~\ref{sec:defects}) rather than
silently using the corrected value without comment.}}
\label{tab:synth-perf}
\end{table}

We report a plain speedup range, $7.0$--$9.3\times$, rather than a
single headline number, because the three independent sessions
summarized in Table~\ref{tab:synth-perf} do not agree to within
measurement noise on this specific circuit family and hardware
configuration. We consider the range itself, and the reasons behind
it (Section~\ref{sec:variance}), more informative than any single
figure drawn from it.

\section{Layout search}
\label{sec:layout}

\subsection{Background and scope}

A companion paper~\cite{vf2cliff} establishes that Qiskit's
\texttt{VF2Layout} pass fails---reporting no solution where a valid
embedding exists---on a specific, reproducible class of saturated
layout instances, and that a single \texttt{rustworkx.vf2\_mapping}
call with a fixed node ordering (\texttt{id\_order=True}) resolves
every tested instance of that class in microseconds. This paper does
not re-derive that finding; it reports the engineering built on top of
it and its independent, end-to-end verification.

PSF-Zero's own layout prototype, \texttt{smart\_vf2\_layout}, predates
that finding and originally used a different strategy: a sequence of
BFS-family node orderings, tried with \texttt{id\_order=True}, falling
back to \texttt{rustworkx}'s own heuristic ordering if all failed. We
found and fixed four defects in this prototype in the course of
incorporating the companion paper's finding, summarized in
Table~\ref{tab:fixes} and detailed in Section~\ref{sec:defects}.

\begin{table}[h]
\centering
\begin{tabular}{p{0.3\textwidth}p{0.4\textwidth}p{0.22\textwidth}}
\toprule
Defect & Fix & Effect \\
\midrule
Natural ordering (matching the companion paper's finding) never
attempted & Added as the first candidate ordering & $56.8\times$ median
speedup on the search step \\
Feasibility guard compared interaction-graph edge count to device
matching size, valid only for the prototype's original circuit family &
Compares the interaction graph's own maximum matching instead & Every
chain-shaped circuit family used in this project reaches the search at
all for the first time \\
Guard recomputed the device's own matching on every call &
Cached per device, keyed on graph structure & $52.4\%$ reduction in
guard cost \\
Ordering candidates computed eagerly regardless of use &
Generator-based, lazy evaluation & Modest additional saving; the larger
residual cost was traced to graph construction, not ordering \\
\bottomrule
\end{tabular}
\caption{Defects found and fixed in \texttt{smart\_vf2\_layout}.}
\label{tab:fixes}
\end{table}

\subsection{Is the feasibility guard worth its cost?}

The repaired guard adds a fixed per-call cost. We measured directly
whether this is worthwhile, using instances specifically constructed to
be infeasible by matching capacity but not by raw qubit count---the
only case in which the guard can add value, since qubit-count
infeasibility is caught by the search itself immediately regardless.

\begin{table}[h]
\centering
\begin{tabular}{lrr}
\toprule
Instance & Guard rejection & Unguarded search \\
\midrule
Matching-imbalanced, small & $0.70$\,ms & $1{,}642.5$\,ms \\
Matching-imbalanced, larger & $1.22$\,ms & $\geq 2{,}000.5$\,ms (budget-limited) \\
Matching-imbalanced, off-by-one & $1.18$\,ms & $\geq 2{,}000.1$\,ms (budget-limited) \\
Control (also infeasible by qubit count) & $1.00$\,ms & $0.30$\,ms \\
\bottomrule
\end{tabular}
\caption{Guard cost versus unguarded search cost on constructed
infeasible instances. Two of four searches did not conclude within the
$2$-second budget used and are reported as lower bounds. The control
row confirms the guard is not simply always cheaper---on an instance
the search itself rejects immediately, the guard's own fixed cost
exceeds it.}
\label{tab:guard}
\end{table}

Excluding the control case (where the guard is not the relevant
mechanism), the guard breaks even against the unguarded search at
roughly one infeasible call in 482, using the smallest, fully-concluded
search time in Table~\ref{tab:guard} as a conservative bound.

\subsection{End-to-end verification}
\label{sec:e2e}

We verified the repaired search through
\texttt{compile\_for\_hardware(layout\_search=True)}, the system's own
public entry point, rather than through \texttt{smart\_vf2\_layout}
alone, on saturated instances at two grid sizes ($6\times7$, $8\times8$)
in two circuit families: the prototype's original disjoint-pair family,
and a chain-shaped family the pre-fix guard rejected outright.

\textbf{Coverage.} All 26 tested instances (matching the companion
paper's own primary instance family) find a valid layout through this
entry point, versus 0 of 26 before the guard fix---since the guard
previously rejected every one before any search ran.

\textbf{Structural validity.} Every routed output was checked directly
against the coupling map's own edge set. Across all measured
configurations, coupling-map violations: $0$.

\textbf{Unitary equivalence.} At $n=6$ (both circuit families, three
arms, three seeds, two repeats: 36 runs), we compared each output's
operator against the input's, correcting for the routing-induced qubit
permutation using Qiskit's own \texttt{TranspileLayout} interface. The
permutation-handling logic itself was independently verified against
hand-constructed test cases before use, since an error there could
produce a false pass or false fail indistinguishable from a genuine
result. Minimum fidelity across all 36 runs: $0.999999999999994$---at
the floating-point noise floor, not near any threshold.

\textbf{Practical benefit.} On the chain-shaped family specifically
(the instances the guard fix newly unlocks), the repaired search
produces a modest but consistent end-to-end improvement over the
default layout path: $1.07\times$ faster at both grid sizes tested,
never worse. We report this without embellishment: it is a real but
small effect, not the multiple-orders-of-magnitude figure that applies
to the isolated search step in Table~\ref{tab:fixes}, because the
search is a small fraction of the full compilation pipeline's own cost.

\section{Generalization}
\label{sec:generalization}

Every result in Section~\ref{sec:layout} used square grids. We tested
the repaired search on two structurally different topologies:
\texttt{brick} (a lattice on which the prototype's own earlier BFS
orderings had previously found nothing on saturated instances) and
heavy-hex (bipartite-imbalanced, so qubit-count saturation is
unattainable; we use saturation relative to the lattice's own maximum
matching instead).

\begin{table}[h]
\centering
\begin{tabular}{lrrr}
\toprule
Topology & $n$ & Repaired search & Prior orderings \\
\midrule
\texttt{brick}, $6\times7$ & 42 & found, $1.25$\,ms & none found, $875$\,ms \\
\texttt{brick}, $8\times8$ & 64 & found, $1.08$\,ms & none found, $996$\,ms \\
heavy-hex, $d{=}3$ & 19 & found, $0.25$\,ms & --- \\
heavy-hex, $d{=}5$ & 57 & found, $0.69$\,ms & --- \\
\bottomrule
\end{tabular}
\caption{Repaired search versus the prototype's own prior ordering
strategies, on topologies neither was developed against.}
\label{tab:generalization}
\end{table}

The \texttt{brick} result is, to us, the most striking evidence that
this is not a square-grid-specific fix: the repaired search succeeds in
about a millisecond on instances where the prototype's entire prior
strategy---every BFS-family ordering it previously used, exhausted---found
nothing at all.

\section{Defects found}
\label{sec:defects}

We report every defect found during this investigation, including two
that reflect on our own prior conclusions rather than only on the
system under test.

\textbf{The four layout-search defects} are summarized in
Table~\ref{tab:fixes} above.

\textbf{A root-cause mischaracterization, corrected.} An early stage of
this investigation attributed part of the guard's inefficiency to a
specific code pattern (\texttt{find\_bit}-based qubit resolution) found
in an intermediate development revision. Re-reading the actual
production codebase directly showed that revision was never deployed;
the pattern in question does not appear in the code this paper
evaluates. We retract the earlier attribution rather than let it stand
alongside the corrected account.

\textbf{A cumulative-speedup figure, corrected.} Table~\ref{tab:synth-perf}
reports the archive session's own speedup as $7.00\times$; our own
earlier project record reported $8.87\times$ for the same session. We
recomputed directly from that session's own raw per-call data while
preparing this paper and found the discrepancy; we do not know whether
the earlier figure was a transcription error or drawn from a different
run than the archived raw data represents, and we do not guess between
these. We use the recomputed, independently-verifiable figure
throughout this paper.

\textbf{A memory-growth question, resolved.} Disabling Python's garbage
collector during repeated \texttt{compile()} calls revealed measurable
memory growth. We traced this to two independent, distinct phenomena:
reference cycles in Qiskit's own circuit objects (not PSF-Zero's own
code), reclaimed cleanly and at a bounded, steady rate under normal
(garbage-collection-enabled) operation; and a separate, one-time
interpreter warm-up cost in generic runtime objects, confirmed by
extending the measurement across multiple checkpoints within a single
process and observing growth reach exactly zero after the first
several hundred calls. Neither is a genuine leak.

\section{An unresolved finding}
\label{sec:variance}

Table~\ref{tab:synth-perf}'s standard-deviation column shows a
$34$--$54\times$ reduction between the archive session and the two
later sessions (confirmed by Levene's test, $p \approx 0$ for both
comparisons against the archive; the two later sessions do not differ
significantly from each other for the \texttt{verify=True} arm,
$p=0.082$). We investigated this directly rather than let it pass
unremarked.

We obtained the actual source code from both sessions and examined the
diff line by line. Two changes present between the sessions were
traced as candidate explanations and ruled out: a change to the gate
synthesizer's own internal caching eviction policy does not
plausibly matter for this benchmark, because the benchmark generates a
fresh, uniquely-seeded random circuit every iteration, so the cache
provides negligible hit-rate benefit under either eviction policy; and
a change eliminating redundant verification work applies only to a
different code path (\texttt{verify="strict"}) not exercised by this
benchmark. The compiled Rust core's own decomposition routine is
unchanged between sessions in every way that could plausibly affect
per-call timing. We found no other candidate in the diff.

We did not investigate two remaining candidate explanations: differences
in the versions of dependency packages installed on the measurement
machine between sessions (not tracked at the time), and machine or
operating-system-level environmental factors. We report this as a
closed but unresolved investigation rather than speculate further.

\section{Related work}

Cartan/KAK decomposition of two-qubit unitaries is
well-established~\cite{kak1,kak2}; our contribution is not the
decomposition itself but a compiled, production-oriented
implementation and its integration into a verified compilation
pipeline. The layout search builds directly on the ordering-sensitivity
finding of our companion paper~\cite{vf2cliff}; SABRE~\cite{sabre}
remains the relevant point of comparison for heuristic routing when no
perfect layout exists.

\section{Conclusion}

PSF-Zero combines an analytic gate synthesizer with a layout search
engineered around a specific, externally-motivated ordering insight,
and we have verified both components more thoroughly than is typical
for a research prototype: exact correctness at small scale, structural
validity and unitary equivalence end-to-end, and generalization beyond
the topology used in development. We have also been direct about what
remains unresolved---an unexplained variance reduction in our own
timing measurements---and about two mistakes in our own prior reporting
that we found and corrected while preparing this paper. We consider
this standard of self-correction a necessary part of the contribution,
not a caveat to it.

\section*{Data and code availability}

All raw measurement data, defect reports, and the complete
chronological verification record, including every correction
discussed in Section~\ref{sec:defects}, are available at
\url{https://github.com/TN-Holdings-LLC/psf-zero}.

\begin{thebibliography}{9}

\bibitem{vf2cliff}
TN Holdings,
\emph{Ordering Sensitivity in Subgraph-Isomorphism Layout Search: A
Characterization of Catastrophic Failure Regions in Quantum Circuit
Transpilation}, 2026.

\bibitem{kak1}
N.~Khaneja and S.~J. Glaser,
\emph{Cartan decomposition of $SU(2^n)$ and control of spin systems},
Chemical Physics, 267(1--3):11--23, 2001.

\bibitem{kak2}
R.~R. Tucci,
\emph{An introduction to Cartan's KAK decomposition for QC
programmers}, arXiv:quant-ph/0507171, 2005.

\bibitem{sabre}
G.~Li, Y.~Ding, and Y.~Xie,
\emph{Tackling the qubit mapping problem for NISQ-era quantum devices},
ASPLOS '19, pp.~1001--1014, 2019.

\end{thebibliography}

\end{document}
